\documentclass[12pt]{article}

\usepackage{amsmath, amssymb}
\usepackage{geometry}
\usepackage{setspace}
\usepackage{booktabs}
\usepackage{graphicx}
\usepackage{xcolor}
\usepackage{mdframed}
\usepackage{enumitem}
\usepackage{hyperref}
\usepackage{tikz}

\geometry{margin=1.2in}
\onehalfspacing

\definecolor{boxblue}{RGB}{220, 235, 250}
\definecolor{boxgreen}{RGB}{220, 245, 220}
\definecolor{boxyellow}{RGB}{255, 248, 220}
\definecolor{boxred}{RGB}{250, 225, 225}

\newmdenv[backgroundcolor=boxblue, linecolor=blue!40, linewidth=1pt,
          innertopmargin=8pt, innerbottommargin=8pt,
          innerleftmargin=10pt, innerrightmargin=10pt]{conceptbox}

\newmdenv[backgroundcolor=boxgreen, linecolor=green!50!black, linewidth=1pt,
          innertopmargin=8pt, innerbottommargin=8pt,
          innerleftmargin=10pt, innerrightmargin=10pt]{intuitionbox}

\newmdenv[backgroundcolor=boxyellow, linecolor=orange!60, linewidth=1pt,
          innertopmargin=8pt, innerbottommargin=8pt,
          innerleftmargin=10pt, innerrightmargin=10pt]{examplebox}

\newmdenv[backgroundcolor=boxred, linecolor=red!40, linewidth=1pt,
          innertopmargin=8pt, innerbottommargin=8pt,
          innerleftmargin=10pt, innerrightmargin=10pt]{warningbox}

\title{\textbf{How Does a Policy Work When Everyone Starts at a Different Time?}\\[6pt]
\large A Gentle Introduction to Bayesian Grouping in Staggered\\
Difference-in-Differences Using the Dirichlet Process}
\author{}
\date{}

\begin{document}

\maketitle
\tableofcontents
\newpage

% =============================================================================
\section{The Problem in Plain English}
% =============================================================================

Imagine a government rolls out a job-training programme across many cities, but
not all at once. Some cities start in January, others in March, and some not
until September. A year later, an economist wants to answer the question:
\textit{Did the programme raise workers' wages?}

This sounds simple, but there is a hidden complication. Cities that started the
programme early might see wage gains that are different from cities that started
late --- perhaps because the early cities had better instructors, or because the
labour market was in a different phase. In other words, \textbf{the treatment
effect might be different for different cities, at different points in time}.

Economists call each city-and-start-time combination a \textbf{cohort}, and the
treatment effect for cohort $g$ at time $t$ is called the
\textbf{Cohort-Average Treatment effect on the Treated}, or \textbf{CATT}
$\tau_{gt}$. With many cohorts and many time periods, there can be dozens of
CATTs to estimate.

The core question this note addresses is:

\begin{conceptbox}
\textbf{Key Question:} Must we treat every CATT as a completely separate number,
or can we group some of them together because they are really measuring the same
underlying effect? And if so, \textit{how do we let the data tell us which
CATTs belong in the same group}?
\end{conceptbox}

% =============================================================================
\section{Two Bad Extremes}
% =============================================================================

Before introducing the solution, it helps to understand the two obvious but
flawed approaches.

\subsection{Extreme 1: Treat Every CATT Separately (Fully Flexible)}

Estimate one number per cohort-time cell. If you have 3 cohorts and 8 time
periods, that is up to 12 separate numbers to estimate.

\begin{itemize}
  \item \textbf{Upside:} You never force effects together that are actually
    different. Each estimate is unbiased.
  \item \textbf{Downside:} Each estimate is noisy. You are using a smaller
    slice of the data for each number. Think of trying to estimate the average
    height of people in a city using only 10 people vs.\ 1,000 people ---
    the estimate with 10 people is unbiased but jumpy.
\end{itemize}

\subsection{Extreme 2: Assume All CATTs Are Identical (Fully Pooled)}

Estimate one single treatment effect for everyone, ignoring when they started.
This is what the classic two-way fixed effects (TWFE) regression does by default.

\begin{itemize}
  \item \textbf{Upside:} Very precise. You use all the data for one number.
  \item \textbf{Downside:} If the early-starting cities really do have a
    different effect than the late-starting cities, you are averaging across
    apples and oranges. The estimate is \textit{biased} --- it does not
    accurately represent the effect for any particular group.
\end{itemize}

\begin{warningbox}
\textbf{The Bias--Variance Tradeoff:} Being flexible avoids bias but gives noisy
estimates (high variance). Pooling gives precise estimates (low variance) but
risks bias if the groups truly differ. We want the best of both worlds.
\end{warningbox}

\subsection{The Middle Ground}

What if some CATTs are genuinely equal and some are not? Suppose cities starting
in January and March both see a \$2/hour wage gain, while cities starting in
September see a \$4/hour gain. Then:
\begin{itemize}
  \item We \textit{should} pool the January and March cohorts (same true effect,
    pooling gains precision without bias).
  \item We \textit{should not} pool them with the September cohort (different
    true effect, pooling would cause bias).
\end{itemize}

The challenge is that we do not know which CATTs are equal --- that is exactly
what we want the data to tell us.

% =============================================================================
\section{What Is a Partition?}
% =============================================================================

A \textbf{partition} is just a way of grouping the CATTs into clusters where
every CATT in the same cluster shares one common treatment effect.

\begin{examplebox}
\textbf{Example.} Suppose we have 4 CATTs labelled $\tau_1, \tau_2, \tau_3,
\tau_4$. Some possible partitions are:
\begin{itemize}[nosep]
  \item $\{\{1,2,3,4\}\}$ --- all four in one group (fully pooled).
  \item $\{\{1\},\{2\},\{3\},\{4\}\}$ --- all separate (fully flexible).
  \item $\{\{1,2\},\{3,4\}\}$ --- two groups of two.
  \item $\{\{1,2,3\},\{4\}\}$ --- first three together, fourth alone.
\end{itemize}
With just 4 CATTs there are already 15 possible partitions. With 12 CATTs
(our simulation) there are over 4 million!
\end{examplebox}

The problem is to find the \textit{right} partition from the data without
checking all 4 million possibilities.

% =============================================================================
\section{The Bayesian Approach: Letting the Data Vote}
% =============================================================================

In a Bayesian approach, instead of picking one partition and declaring it
correct, we assign a \textbf{probability} to every partition. Partitions that
fit the data well get high probability; partitions that fit poorly get low
probability. Our final answer averages over all partitions, weighted by their
probabilities.

\begin{conceptbox}
\textbf{Bayesian Logic in One Sentence:}\\
\textit{Start with a prior belief about which partitions are plausible, update
that belief using the data, and end up with a posterior belief that reflects
both.}
\end{conceptbox}

This requires two ingredients:
\begin{enumerate}
  \item \textbf{A prior} --- our initial guess about the grouping structure
    \textit{before} looking at the data.
  \item \textbf{A likelihood} --- how probable the data are \textit{given} a
    particular grouping.
\end{enumerate}

For the likelihood, we use our regression model: given a grouping, the data
follow a normal distribution around the group's shared treatment effect.

For the prior, we use something called the \textbf{Dirichlet Process}.

% =============================================================================
\section{The Dirichlet Process: A Prior Over Groupings}
% =============================================================================

\subsection{The Intuition}

The Dirichlet Process (DP) is a probability distribution whose outcomes are
themselves probability distributions. That sounds abstract, so here is the key
intuition for our setting:

\begin{intuitionbox}
\textbf{What the DP does:} It provides a prior that says ``CATTs that are
similar in their true effects will tend to be grouped together, and the number
of groups is determined by the data.'' You do not need to specify in advance
how many groups there are.
\end{intuitionbox}

The DP has two inputs:
\begin{itemize}
  \item $\alpha$ (the \textbf{concentration parameter}): controls how many
    groups you expect. Small $\alpha$ means few groups (like the pooled
    estimator); large $\alpha$ means many groups (like the flexible estimator).
  \item $G_0$ (the \textbf{base measure}): your prior belief about what a
    treatment effect looks like for a brand-new, ungrouped CATT. We use a
    normal distribution centred at zero.
\end{itemize}

\subsection{The Chinese Restaurant Process: A Story}

The easiest way to understand the DP prior on groupings is through a
story called the \textbf{Chinese Restaurant Process (CRP)}.

Imagine a Chinese restaurant with infinitely many tables. CATTs arrive one by
one and choose where to sit:

\begin{examplebox}
\textbf{Chinese Restaurant Story}
\begin{enumerate}
  \item CATT 1 arrives and always sits at Table 1 (the first table).
  \item CATT 2 arrives. It can:
    \begin{itemize}[nosep]
      \item Sit with CATT 1 at Table 1, with probability $\frac{1}{1+\alpha}$.
      \item Start a new Table 2, with probability $\frac{\alpha}{1+\alpha}$.
    \end{itemize}
  \item CATT 3 arrives. It sits at a table with probability proportional to
    how many CATTs are already there, or starts a new table with probability
    proportional to $\alpha$.
  \item \ldots and so on.
\end{enumerate}
The rule is: \textbf{popular tables get more popular} (``rich get richer''),
but there is always a chance of opening a new table.
\end{examplebox}

The seating arrangement at the end --- which CATT sits at which table --- is
a partition of the CATTs into groups. Each table corresponds to one group, and
all CATTs at the same table share one treatment effect.

\textbf{How $\alpha$ controls grouping:}
\begin{center}
\begin{tabular}{cll}
\toprule
$\alpha$ & Tendency & Analogy \\
\midrule
Small (e.g., 0.1) & Few tables, most CATTs together & Strong prior for pooling \\
Medium (e.g., 1) & A few tables, data-driven & Weakly informative \\
Large (e.g., 10) & Many tables, CATTs spread out & Prior for flexibility \\
\bottomrule
\end{tabular}
\end{center}

\subsection{The Mathematical Statement}

For $K$ CATTs, the CRP assigns cluster labels
$z_1, z_2, \ldots, z_K$ sequentially. For CATT $k$:
\begin{equation}
  \Pr(z_k = c \mid z_1, \ldots, z_{k-1}) =
  \begin{cases}
    \dfrac{n_c}{k - 1 + \alpha} & \text{if table } c \text{ already has } n_c \text{ CATTs},\\[8pt]
    \dfrac{\alpha}{k - 1 + \alpha} & \text{if } c \text{ is a new, empty table.}
  \end{cases}
\end{equation}

Once all CATTs are seated, each table $l$ draws its shared treatment effect
from the base measure:
\begin{equation}
  \phi_l \sim G_0 = \mathcal{N}(\mu_0, \sigma_0^2).
\end{equation}
All CATTs at table $l$ then have $\theta_k = \phi_l$.

% =============================================================================
\section{Putting It Together: The Full Model}
% =============================================================================

Here is the complete Bayesian model in plain steps:

\begin{conceptbox}
\textbf{The Full Model (top to bottom)}

\medskip
\textbf{Step 1 --- Assign CATTs to groups (CRP prior):}
\[
  z_k \mid z_1,\ldots,z_{k-1} \sim \text{CRP}(\alpha)
\]
This is the prior. It says: CATTs that sit at the same table share one effect.

\medskip
\textbf{Step 2 --- Draw a treatment effect for each table:}
\[
  \phi_l \sim \mathcal{N}(\mu_0,\, \sigma_0^2)
\]
Each table's effect is drawn from a normal distribution.

\medskip
\textbf{Step 3 --- Generate the data (likelihood):}
\[
  y_{it} = \alpha_i + \lambda_t + \sum_k \phi_{z_k} \cdot D_{kit} + \varepsilon_{it},
  \qquad \varepsilon_{it} \sim \mathcal{N}(0, \sigma^2)
\]
The outcome follows the TWFE model, where CATTs in the same group
share the same coefficient.
\end{conceptbox}

After collecting the data, we want to go \textit{backwards}: given the
outcomes we observed, what is the probability of each possible grouping?
That is the \textbf{posterior}, computed via Bayes' rule:
\begin{equation}
  \underbrace{\Pr(\text{grouping} \mid \text{data})}_{\text{posterior}}
  \;\propto\;
  \underbrace{p(\text{data} \mid \text{grouping})}_{\text{likelihood}}
  \;\times\;
  \underbrace{\Pr(\text{grouping})}_{\text{CRP prior}}.
\end{equation}

% =============================================================================
\section{How Do We Compute the Posterior? The Gibbs Sampler}
% =============================================================================

We cannot enumerate all 4 million+ partitions. Instead, we use an algorithm
called a \textbf{Gibbs sampler} to \textit{draw random samples} from the
posterior. Think of it like exploring a city by taking many random walks ---
you will spend more time in neighbourhoods that are large and interesting
(high-probability partitions) and less time in dead ends (low-probability ones).

\subsection{The Algorithm in Plain English}

\begin{examplebox}
\textbf{One Gibbs Iteration}

Repeat the following for each CATT $k = 1, 2, \ldots, K$:

\begin{enumerate}
  \item \textbf{Kick CATT $k$ out of its current group.} (Temporarily ignore
    where it sits.)

  \item \textbf{For each possible new seat} (every existing group, plus a brand
    new one):
    \begin{itemize}[nosep]
      \item Compute: how much does the data fit if I put CATT $k$ here?
      \item Compute: how popular is this table already? (CRP weight)
      \item Multiply these two numbers together.
    \end{itemize}

  \item \textbf{Randomly assign CATT $k$} to one of the seats, with probability
    proportional to the numbers computed above.

  \item After all $K$ CATTs have been updated: \textbf{draw a new treatment
    effect $\phi_l$} for each group from its posterior (using the data from
    that group only).
\end{enumerate}
\end{examplebox}

After many iterations, the sequence of groupings visited by the algorithm
forms a representative sample from the posterior. We then read off the answers
from this sample.

\subsection{Why Does the Data Fit Calculation Work?}

The key step is: \textit{how much does the data fit if CATT $k$ is in group $c$?}

The answer uses our normal regression model. When CATT $k$ joins group $c$,
the data from CATT $k$ are now explained by the same coefficient $\phi_c$ as
all other CATTs in group $c$. If CATT $k$'s data look like they have a similar
effect to the rest of group $c$, the data fit well (high probability). If
CATT $k$ looks very different, the fit is poor (low probability).

Formally, we integrate out the group effect $\phi_c$ to get the
\textbf{marginal likelihood}:
\begin{equation}
  p(\tilde{y} \mid \text{CATT } k \text{ in group } c)
  = \int p(\tilde{y} \mid \phi_c,\, \text{CATT } k \text{ in group } c)\;
          p(\phi_c)\; d\phi_c,
\end{equation}
where $\tilde{y}$ is the within-transformed outcome (after removing unit and
time fixed effects). Because the prior and likelihood are both normal, this
integral has a closed-form solution --- no numerical approximation needed.

% =============================================================================
\section{What Do We Get Out?}
% =============================================================================

After running the Gibbs sampler for $S$ iterations, we have $S$ draws of the
grouping structure and $S$ draws of the group effects. From these we extract:

\subsection{Posterior Mean of Each CATT}

For each CATT $k$, the posterior mean treatment effect is:
\begin{equation}
  \hat{\theta}_k = \frac{1}{S}\sum_{s=1}^{S} \phi_{z_k^{(s)}}^{(s)},
\end{equation}
where $z_k^{(s)}$ is the group that CATT $k$ belongs to in draw $s$, and
$\phi_{z_k^{(s)}}^{(s)}$ is the corresponding group effect. This automatically
averages over uncertainty in the grouping.

\subsection{Credible Intervals}

A 95\% credible interval for $\theta_k$ is simply the 2.5th and 97.5th
percentiles of $\{\phi_{z_k^{(s)}}^{(s)}\}_{s=1}^{S}$. Unlike classical
confidence intervals, this has the direct interpretation: \textit{``There is a
95\% probability that the true CATT lies in this range.''}

\subsection{Probability That Two CATTs Are Equal}

We can directly answer: \textit{What is the probability that cohort 3 and
cohort 5 have the same treatment effect?}
\begin{equation}
  \hat{\pi}_{jk} = \Pr(\theta_j = \theta_k \mid \text{data})
  \approx \frac{1}{S}\sum_{s=1}^{S} \mathbf{1}\{z_j^{(s)} = z_k^{(s)}\}.
\end{equation}
This is just the fraction of draws in which CATTs $j$ and $k$ end up in the
same group. If $\hat{\pi}_{jk} \approx 1$, the data strongly suggest they
should be pooled. If $\hat{\pi}_{jk} \approx 0$, the data say keep them
separate.

\subsection{Posterior Distribution Over Number of Groups}

We can also plot how many groups appeared in each draw. This tells us not just
``there are 2 groups'' but rather ``there is a 90\% chance of 2 groups and a
10\% chance of 3 groups.''

% =============================================================================
\section{A Numerical Example}
% =============================================================================

To make this concrete, here are the results from running 100 Gibbs iterations
on our simulated data. Recall the true setup:
\begin{itemize}
  \item \textbf{Group A} (10 CATTs): cohorts entering at $t=3$ and $t=5$,
    true effect $\tau = 2$.
  \item \textbf{Group B} (2 CATTs): cohort entering at $t=7$,
    true effect $\tau = 4$.
\end{itemize}

\begin{center}
\begin{tabular}{lcc}
\toprule
Comparison & $\hat\pi$ (posterior co-clustering prob.) & Verdict \\
\midrule
Two CATTs, both in Group A ($\tau=2$) & 0.954 & Almost always pooled \\
Two CATTs, both in Group B ($\tau=4$) & 0.990 & Almost always pooled \\
One CATT from A, one from B           & 0.000 & Never incorrectly merged \\
\bottomrule
\end{tabular}
\end{center}

The posterior mean number of groups is 2.1, closely matching the true value of
2. The posterior densities for each CATT concentrate tightly around the true
effect values (2 or 4), with negligible mass elsewhere.

% =============================================================================
\section{How Is This Different From the $\ell_0$ Greedy Approach?}
% =============================================================================

Both the DP mixture estimator and the $\ell_0$ greedy estimator try to find
the right grouping of CATTs. The table below summarises the key differences:

\begin{center}
\begin{tabular}{p{4.5cm} p{4.5cm} p{4.5cm}}
\toprule
 & \textbf{$\ell_0$ Greedy} & \textbf{DP Mixture (Bayesian)} \\
\midrule
\textbf{Output} & One best-guess partition & Full distribution over partitions \\
\addlinespace
\textbf{Uncertainty} & None --- hard decision & Full posterior uncertainty \\
\addlinespace
\textbf{Tuning parameter} & $L$ (penalty per CATT pair) & $\alpha$ (DP concentration) \\
\addlinespace
\textbf{Number of groups} & Fixed by $L$ & Determined by data \\
\addlinespace
\textbf{Computation} & Very fast ($O(K^3)$ merges) & Slower (MCMC iterations) \\
\addlinespace
\textbf{Credible intervals} & Not available & Naturally produced \\
\addlinespace
\textbf{Conceptual connection} & MAP estimate & Full posterior \\
\bottomrule
\end{tabular}
\end{center}

\begin{intuitionbox}
\textbf{Analogy:} The $\ell_0$ greedy method is like a detective who looks at
all the evidence and announces a single verdict: ``These CATTs are equal, those
are not.'' The DP Bayesian method is like a jury that assigns a probability to
each possible verdict and can say: ``We are 95\% confident these two CATTs
belong in the same group.''
\end{intuitionbox}

% =============================================================================
\section{Limitations and Things to Keep in Mind}
% =============================================================================

\begin{enumerate}

  \item \textbf{Convergence.} The Gibbs sampler needs enough iterations to
    explore the posterior thoroughly. With only 100 iterations on a problem
    with 12 CATTs and a clear true partition, it works well. For larger or
    noisier problems, more iterations (and a burn-in period to discard early
    draws) may be needed.

  \item \textbf{Choosing $\alpha$.} The concentration parameter $\alpha$
    influences the prior probability of how many groups exist. In practice,
    you can put a prior on $\alpha$ itself and let the data update it, rather
    than fixing it by hand.

  \item \textbf{Choosing $G_0$.} The base measure $G_0 = \mathcal{N}(\mu_0,
    \sigma_0^2)$ encodes your prior belief about the size of treatment effects
    for any new group. Setting $\sigma_0^2$ very large (diffuse) makes the
    prior non-informative and lets the data dominate.

  \item \textbf{Pre-trends.} Like all DiD methods, the Bayesian DP approach
    requires that in the absence of treatment, all cohorts would have followed
    parallel trends. This assumption cannot be tested with post-treatment data.

  \item \textbf{Not a free lunch.} The Bayesian approach quantifies uncertainty
    about the grouping structure, but it does not eliminate it. If the true
    treatment effects are very close together (e.g., 2.0 vs.\ 2.1), even a
    Bayesian model with lots of data may not reliably separate them.

\end{enumerate}

% =============================================================================
\section{Summary}
% =============================================================================

Here is the complete story in five bullet points:

\begin{itemize}

  \item In staggered DiD, each cohort-time cell has its own treatment effect
    (CATT). Estimating all of them separately is unbiased but noisy; pooling
    them all is precise but biased.

  \item A smarter approach: group CATTs that truly share the same effect, and
    let the data determine the grouping. This is the \textit{model specification}
    problem.

  \item The Dirichlet Process provides a prior over all possible groupings,
    governed by one intuitive parameter $\alpha$. The Chinese Restaurant Process
    story explains how groups form: popular groups attract more members, but new
    groups can always be opened.

  \item A Gibbs sampler explores the posterior over groupings by updating each
    CATT's group assignment one at a time, balancing how well the data fit with
    how many CATTs are already in each group.

  \item The output is not just one estimated partition but a full probability
    distribution: which CATTs likely share an effect, how many groups probably
    exist, and how certain we are about each CATT's treatment effect.

\end{itemize}

\bigskip
\noindent\textbf{Further Reading}
\begin{itemize}[nosep]
  \item Callaway \& Sant'Anna (2021) --- how to estimate CATTs in staggered DiD.
  \item Neal (2000) --- the collapsed Gibbs sampler for DP mixture models.
  \item M\"{u}ller \& Quintana (2004) --- a friendly introduction to Bayesian
    nonparametric methods for applied researchers.
\end{itemize}

\end{document}
